% Options for packages loaded elsewhere
\PassOptionsToPackage{unicode}{hyperref}
\PassOptionsToPackage{hyphens}{url}
\documentclass[
]{article}
\usepackage{xcolor}
\usepackage{amsmath,amssymb}
\setcounter{secnumdepth}{-\maxdimen} % remove section numbering
\usepackage{iftex}
\ifPDFTeX
  \usepackage[T1]{fontenc}
  \usepackage[utf8]{inputenc}
  \usepackage{textcomp} % provide euro and other symbols
\else % if luatex or xetex
  \usepackage{unicode-math} % this also loads fontspec
  \defaultfontfeatures{Scale=MatchLowercase}
  \defaultfontfeatures[\rmfamily]{Ligatures=TeX,Scale=1}
\fi
\usepackage{lmodern}
\ifPDFTeX\else
  % xetex/luatex font selection
\fi
% Use upquote if available, for straight quotes in verbatim environments
\IfFileExists{upquote.sty}{\usepackage{upquote}}{}
\IfFileExists{microtype.sty}{% use microtype if available
  \usepackage[]{microtype}
  \UseMicrotypeSet[protrusion]{basicmath} % disable protrusion for tt fonts
}{}
\makeatletter
\@ifundefined{KOMAClassName}{% if non-KOMA class
  \IfFileExists{parskip.sty}{%
    \usepackage{parskip}
  }{% else
    \setlength{\parindent}{0pt}
    \setlength{\parskip}{6pt plus 2pt minus 1pt}}
}{% if KOMA class
  \KOMAoptions{parskip=half}}
\makeatother
\usepackage{graphicx}
\makeatletter
\newsavebox\pandoc@box
\newcommand*\pandocbounded[1]{% scales image to fit in text height/width
  \sbox\pandoc@box{#1}%
  \Gscale@div\@tempa{\textheight}{\dimexpr\ht\pandoc@box+\dp\pandoc@box\relax}%
  \Gscale@div\@tempb{\linewidth}{\wd\pandoc@box}%
  \ifdim\@tempb\p@<\@tempa\p@\let\@tempa\@tempb\fi% select the smaller of both
  \ifdim\@tempa\p@<\p@\scalebox{\@tempa}{\usebox\pandoc@box}%
  \else\usebox{\pandoc@box}%
  \fi%
}
% Set default figure placement to htbp
\def\fps@figure{htbp}
\makeatother
\setlength{\emergencystretch}{3em} % prevent overfull lines
\providecommand{\tightlist}{%
  \setlength{\itemsep}{0pt}\setlength{\parskip}{0pt}}
% ---------------------------------------------------------------------------
% Unicode -> LaTeX mappings so these documents compile and render correctly
% under xelatex (and pdflatex) without depending on a particular system font.
% Every non-ASCII glyph used in the course text is mapped explicitly; this is
% needed because the documents use mathematical symbols inside running text
% and inside inline-code spans, where unicode-math would otherwise treat them
% as math-only and abort.
% ---------------------------------------------------------------------------
\usepackage{amssymb}
\usepackage{newunicodechar}

% --- Greek (lower case) ---
\newunicodechar{α}{\ensuremath{\alpha}}
\newunicodechar{β}{\ensuremath{\beta}}
\newunicodechar{γ}{\ensuremath{\gamma}}
\newunicodechar{δ}{\ensuremath{\delta}}
\newunicodechar{θ}{\ensuremath{\theta}}
\newunicodechar{κ}{\ensuremath{\kappa}}
\newunicodechar{λ}{\ensuremath{\lambda}}
\newunicodechar{μ}{\ensuremath{\mu}}
\newunicodechar{ν}{\ensuremath{\nu}}
\newunicodechar{π}{\ensuremath{\pi}}
\newunicodechar{ρ}{\ensuremath{\rho}}
\newunicodechar{ς}{\ensuremath{\varsigma}}
\newunicodechar{σ}{\ensuremath{\sigma}}
\newunicodechar{τ}{\ensuremath{\tau}}
\newunicodechar{φ}{\ensuremath{\varphi}}
\newunicodechar{χ}{\ensuremath{\chi}}
% --- Greek (capitals) ---
\newunicodechar{Γ}{\ensuremath{\Gamma}}
\newunicodechar{Δ}{\ensuremath{\Delta}}
\newunicodechar{Λ}{\ensuremath{\Lambda}}
\newunicodechar{Σ}{\ensuremath{\Sigma}}

% --- Math symbols ---
\newunicodechar{ℝ}{\ensuremath{\mathbb{R}}}
\newunicodechar{∂}{\ensuremath{\partial}}
\newunicodechar{∈}{\ensuremath{\in}}
\newunicodechar{∫}{\ensuremath{\int}}
\newunicodechar{√}{\ensuremath{\surd}}
\newunicodechar{≈}{\ensuremath{\approx}}
\newunicodechar{≤}{\ensuremath{\leq}}
\newunicodechar{≥}{\ensuremath{\geq}}
\newunicodechar{±}{\ensuremath{\pm}}
\newunicodechar{×}{\ensuremath{\times}}
\newunicodechar{−}{\ensuremath{-}}
\newunicodechar{·}{\ensuremath{\cdot}}
\newunicodechar{½}{\textonehalf}
\newunicodechar{→}{\ensuremath{\rightarrow}}
\newunicodechar{⟶}{\ensuremath{\longrightarrow}}

% --- Superscripts / modifier letters ---
\newunicodechar{²}{\ensuremath{^{2}}}
\newunicodechar{³}{\ensuremath{^{3}}}
\newunicodechar{⁺}{\ensuremath{^{+}}}
\newunicodechar{⁻}{\ensuremath{^{-}}}
\newunicodechar{ᵀ}{\ensuremath{^{\mathsf{T}}}}
\newunicodechar{ᵏ}{\ensuremath{^{k}}}

% --- Punctuation / misc text symbols ---
\newunicodechar{…}{\textellipsis}
\newunicodechar{§}{\S}
\newunicodechar{⚑}{\textbf{[!]}}

% --- Accented Latin (author names) ---
\newunicodechar{Á}{\'A}
\newunicodechar{Ø}{\O}
\newunicodechar{ä}{\"a}
\newunicodechar{é}{\'e}
\newunicodechar{ó}{\'o}
\newunicodechar{ô}{\^o}
\usepackage{bookmark}
\IfFileExists{xurl.sty}{\usepackage{xurl}}{} % add URL line breaks if available
\urlstyle{same}
\hypersetup{
  pdftitle={Problem Sets},
  pdfauthor={North Carolina State University · Financial Mathematics},
  hidelinks,
  pdfcreator={LaTeX via pandoc}}

\title{Problem Sets}
\usepackage{etoolbox}
\makeatletter
\providecommand{\subtitle}[1]{% add subtitle to \maketitle
  \apptocmd{\@title}{\par {\large #1 \par}}{}{}
}
\makeatother
\subtitle{Optimal Market Making for Cryptocurrency Options --- The
Avellaneda--Stoikov Framework and its Extension to Options}
\author{North Carolina State University · Financial Mathematics}
\date{Summer 2026}

\begin{document}
\maketitle

\includegraphics[width=2.6in,height=\textheight,keepaspectratio]{assets/ncsu_logo.png}

\section{Problem Sets}\label{problem-sets}

\textbf{Optimal Market Making for Cryptocurrency Options}

\emph{The Avellaneda--Stoikov Framework and its Extension to Options}

North Carolina State University --- Summer 2026

\begin{center}\rule{0.5\linewidth}{0.5pt}\end{center}

\subsection{Instructions}\label{instructions}

These are the weekly programming assignments for Weeks 2--10 (Week 1 was
the introduction week and had no graded set).

\begin{itemize}
\tightlist
\item
  \textbf{Problem sets are written in Python}, using \textbf{PyTorch} as
  the numerical engine. Use PyTorch autograd for all derivatives
  (Greeks, Jacobians, gradients) --- you are \textbf{not} asked to
  hand-derive or hand-code algorithmic differentiation.
\item
  \textbf{Week 7 is the exception}: it is written in \textbf{C++} and
  called from Python through a \textbf{\texttt{pybind11}} binding. A
  \texttt{CMakeLists.txt} scaffold is provided.
\item
  Each problem set ships with a starter repository: data loaders, a test
  harness (\texttt{pytest}), and (Week 7) the \texttt{pybind11}/CMake
  scaffold.
\item
  Emit numerical results as CSV and figures as PNG for inclusion in your
  write-up. Where a problem asks for a comment, print it from your code
  as a short summary of the actual numbers.
\item
  Submissions are due Sunday 23:59 Eastern. Set seeds for
  reproducibility.
\end{itemize}

Worked solutions are distributed separately in the \textbf{Solutions}
booklet after each deadline.

\begin{center}\rule{0.5\linewidth}{0.5pt}\end{center}

\subsection{\texorpdfstring{Problem Set 2 --- Foundations of Market
Microstructure
\emph{(Python)}}{Problem Set 2 --- Foundations of Market Microstructure (Python)}}\label{problem-set-2-foundations-of-market-microstructure-python}

\begin{enumerate}
\def\labelenumi{\arabic{enumi}.}
\tightlist
\item
  \textbf{Order-book reconstruction.} Implement a Python class
  \texttt{L2Book} that consumes a stream of events (\texttt{new},
  \texttt{cancel}, \texttt{modify}, \texttt{trade}) and maintains the L2
  state in a price-keyed structure (e.g.~two \texttt{SortedDict}s).
  Apply it to one hour of Coincall BTC-PERP data and emit a CSV of
  mid-price and quoted spread over time. Remember to remove price levels
  that become empty after a cancel.
\item
  \textbf{Roll estimator.} Compute the Roll effective-spread estimator
  \texttt{2*sqrt(-cov(Deltap\_t,\ Deltap\_\{t-1\}))} on the same sample
  and compare it to the directly observed bid-ask spread. Print both and
  their discrepancy.
\item
  \textbf{Glosten--Milgrom quote-setter.} Implement the rational
  quote-setter as a PyTorch belief update: given an informed-trade
  probability and a prior over the asset value, compute bid and ask from
  the Bayesian posterior after a buy or a sell. Unit-test against a
  hand-computed one-trade sequence.
\item
  \textbf{Glosten--Milgrom simulation.} Build a Monte Carlo simulator
  reusing problem 3 and emit a CSV of the public belief converging to
  the true value, averaged over many runs.
\item
  \textbf{Toxicity.} Compute realized spread and price impact over a
  5-minute decomposition window on one hour of Coincall data; print a
  summary relating the numbers to order-flow toxicity.
\end{enumerate}

\subsection{\texorpdfstring{Problem Set 3 --- The Avellaneda--Stoikov
Model
\emph{(Python/PyTorch)}}{Problem Set 3 --- The Avellaneda--Stoikov Model (Python/PyTorch)}}\label{problem-set-3-the-avellanedastoikov-model-pythonpytorch}

\begin{enumerate}
\def\labelenumi{\arabic{enumi}.}
\tightlist
\item
  \textbf{Reduced ODE solver.} Implement a PyTorch backward integrator
  for the reduced system in \texttt{theta(t,\ q)} arising from the
  ansatz \texttt{V\ =\ -exp(-gamma(x+qS))*theta(t,q)}. Emit
  \texttt{theta} on a \texttt{(t,\ q)} grid.
\item
  \textbf{Closed-form quotes.} Implement the asymptotic reservation
  price \texttt{r\ =\ S\ -\ qgammasigma\^{}2(T-t)} and half-spread
  \texttt{delta\ =\ (1/gamma)*ln(1+gamma/kappa)\ +\ (gammasigma\^{}2/2)(T-t)}
  as a reusable function, and \texttt{pytest} it against the solver from
  problem 1 where they should agree.
\item
  \textbf{Quote engine + simulation.} Simulate one trading day with
  \texttt{sigma\ =\ 2\ USD/sqrts}, \texttt{kappa\ =\ 1.5},
  \texttt{A\ =\ 140}. Emit the P\&L distribution of the optimal policy
  alongside a symmetric, inventory-ignoring baseline.
\item
  \textbf{γ sensitivity.} Sweep \texttt{gamma} over three orders of
  magnitude and emit a CSV of inventory variance and Sharpe
  vs.~\texttt{gamma}; identify in code the most effective regime.
\item
  \textbf{Price impact.} Extend the simulator so the mid reacts linearly
  to your own fills; run three impact levels and report how each
  degrades the optimal policy's P\&L (independence-assumption
  violation).
\end{enumerate}

\subsection{\texorpdfstring{Problem Set 4 --- Stochastic Optimal Control
\emph{(Python/PyTorch)}}{Problem Set 4 --- Stochastic Optimal Control (Python/PyTorch)}}\label{problem-set-4-stochastic-optimal-control-pythonpytorch}

\begin{enumerate}
\def\labelenumi{\arabic{enumi}.}
\tightlist
\item
  \textbf{Merton CRRA.} Solve the Merton problem under CRRA by
  discretizing the HJB (value iteration or backward ODE); confirm
  numerically the optimal portfolio \emph{share} is constant in wealth.
\item
  \textbf{Merton CARA.} Repeat under CARA; confirm the optimal
  \emph{dollar amount} is constant in wealth; emit a side-by-side CSV
  against problem 1.
\item
  \textbf{Almgren--Chriss.} Solve continuous-time optimal execution
  under quadratic transient impact; verify it matches the discrete-time
  result in the appropriate limit.
\item
  \textbf{HJB-residual verification.} Write a harness that evaluates the
  HJB residual for the Avellaneda--Stoikov value function (use autograd
  for the derivatives). Show the residual is \textasciitilde machine
  zero for the true solution, then perturb it and show the residual
  diverges.
\item
  \textbf{Obstacle problem.} Solve the obstacle problem (projected SOR
  for the variational inequality) and compare to the unconstrained
  classical solution, demonstrating where a viscosity treatment is
  required.
\end{enumerate}

\subsection{\texorpdfstring{Problem Set 5 --- Avellaneda--Stoikov Engine
\emph{(Python/PyTorch --- Milestone
1)}}{Problem Set 5 --- Avellaneda--Stoikov Engine (Python/PyTorch --- Milestone 1)}}\label{problem-set-5-avellanedastoikov-engine-pythonpytorch-milestone-1}

\begin{enumerate}
\def\labelenumi{\arabic{enumi}.}
\tightlist
\item
  \textbf{Vectorized quote path.} Implement the quote computation
  vectorized over the inventory grid in PyTorch; benchmark it and report
  the latency distribution (mean, median, p95, p99).
\item
  \textbf{Intensity calibration.} Calibrate
  \texttt{Lambda\^{}+,\ Lambda\^{}-} from one week of Coincall fill data
  by maximum likelihood (minimize the negative log-likelihood with
  \texttt{torch.optim}); report fitted \texttt{A}, \texttt{kappa}, and a
  χ² goodness-of-fit.
\item
  \textbf{Paper reproduction.} Reproduce the inventory/P\&L behaviour of
  Avellaneda--Stoikov (2008) on simulated data; emit the figures as CSV;
  print a summary of any discrepancies.
\item
  \textbf{One-month backtest.} Build the event-driven harness (replay,
  fill simulator, P\&L attribution) and run a one-month BTC-PERP
  backtest; report mean/vol of P\&L, Sharpe, max drawdown, mean and std
  of inventory, vs.~a symmetric baseline.
\item
  \textbf{Tests and reproducibility.} Provide a \texttt{pytest} suite
  over the quote computation and cash/inventory bookkeeping; demonstrate
  reproducible backtests under a fixed seed.
\end{enumerate}

\emph{Deliverable: Milestone 1 (see syllabus §2 for acceptance
criteria).}

\subsection{\texorpdfstring{Problem Set 6 --- Volatility Surface and
Calibration
\emph{(Python/PyTorch)}}{Problem Set 6 --- Volatility Surface and Calibration (Python/PyTorch)}}\label{problem-set-6-volatility-surface-and-calibration-pythonpytorch}

\begin{enumerate}
\def\labelenumi{\arabic{enumi}.}
\tightlist
\item
  \textbf{BS Greeks via autograd.} Implement Black--Scholes pricing in
  PyTorch and obtain delta, gamma, vega, theta, vanna, volga from
  autograd. Verify vanna/volga against finite differences; report the
  max relative error across moneyness/maturity.
\item
  \textbf{SVI + arbitrage checks.} Implement raw SVI and the
  Gatheral--Jacquier arbitrage-free constraint checks; apply to a
  deliberately broken parameterization and have the code report which
  constraints fail.
\item
  \textbf{SVI calibration.} Fit SVI to one Coincall BTC chain snapshot
  (\texttt{torch.optim}, autograd Jacobians); report parameters,
  per-strike residuals, and RMS error in vol units.
\item
  \textbf{SABR calibration.} Fit SABR to the same snapshot; emit a
  quantitative comparison against SVI and a printed note on
  strengths/weaknesses for crypto.
\item
  \textbf{Stability study.} Run daily SVI calibration over one month;
  emit the parameter time series; implement a smoothing scheme (seed
  each day from the previous fit) and report its effect.
\end{enumerate}

\subsection{\texorpdfstring{Problem Set 7 --- Fast Computation
\emph{(C++ with \texttt{pybind11}, called from Python --- Milestone
2)}}{Problem Set 7 --- Fast Computation (C++ with pybind11, called from Python --- Milestone 2)}}\label{problem-set-7-fast-computation-c-with-pybind11-called-from-python-milestone-2}

A \texttt{CMakeLists.txt} and \texttt{pybind11} scaffold are provided.
Build the module
(\texttt{cmake\ -B\ build\ \&\&\ cmake\ -\/-build\ build}) and import it
from Python as \texttt{fastmm}.

\begin{enumerate}
\def\labelenumi{\arabic{enumi}.}
\tightlist
\item
  \textbf{COS / Carr--Madan pricer (C++).} Implement the Fang--Oosterlee
  COS pricer and the Carr--Madan FFT pricer for Heston European calls in
  C++ (use the provided FFT dependency). Bind both to Python. Verify
  against a PyTorch Monte Carlo benchmark and report COS convergence
  vs.~\texttt{N}.
\item
  \textbf{Robust implied volatility (C++).} Implement an implied-vol
  inversion with a rational initial guess and a safeguarded
  Newton/Halley step; bind it. Show convergence across the surface,
  including deep wings and short expiries.
\item
  \textbf{Greeks, validated by autograd.} Expose prices from C++ and
  compute the Greek vector. Validate every Greek against a
  \textbf{PyTorch-autograd} reference implementation of the same pricer;
  report max relative error. (Do \textbf{not} hand-code adjoint AD in
  C++ --- autograd is the reference.)
\item
  \textbf{Vectorized surface revaluation (C++).} Revalue the full
  Coincall BTC surface (300--500 contracts) with the Greek vector;
  release the GIL around the C++ loop; report per-option latency (mean,
  median, p95, p99) and the full-surface time (target \textless{} 5 ms).
\item
  \textbf{The binder, explained.} In a short \texttt{BINDING.md},
  explain your \texttt{pybind11} module to a student new to it: how
  tensors cross the boundary without copying (buffer protocol /
  \texttt{py::array\_t}), why and where you release the GIL, and how the
  CMake build finds Python and \texttt{pybind11}.
\end{enumerate}

\emph{Deliverable: Milestone 2 (see syllabus §2 for acceptance
criteria).}

\subsection{\texorpdfstring{Problem Set 8 --- Options Market Making and
Hedging
\emph{(Python/PyTorch)}}{Problem Set 8 --- Options Market Making and Hedging (Python/PyTorch)}}\label{problem-set-8-options-market-making-and-hedging-pythonpytorch}

\begin{enumerate}
\def\labelenumi{\arabic{enumi}.}
\tightlist
\item
  \textbf{Multi-Greek inventory.} Implement inventory aggregation and
  the inventory dynamics under fills with a vector
  \texttt{q\ \ in\ \ R\^{}k}; unit-test the jump structure a fill
  induces.
\item
  \textbf{El Aoud--Abergel solver.} Solve the reduced multi-Greek
  problem under a quadratic penalty in PyTorch; verify in code it
  collapses to the Avellaneda--Stoikov quotes when one Greek and one
  option are tracked.
\item
  \textbf{Penalty-matrix calibration.} Calibrate the quadratic penalty
  matrix from the empirical covariance of delta/gamma/vega over one week
  of Coincall data; emit the matrix with each entry justified by the
  data.
\item
  \textbf{Gamma--theta--variance identity.} Simulate a continuously
  delta-hedged short call (autograd for the hedge ratio) and verify
  numerically the P\&L equals
  \texttt{integral\ \ 0.5*Gamma*(realized\ -\ implied\ variance)}.
\item
  \textbf{Threshold hedger + engine.} Implement a delta-hedger
  triggering when
  \texttt{\textbar{}Delta\textbar{}\ \textgreater{}\ 0.1\ BTC},
  integrate it with the Week-5 engine on one week of data, and report
  hedging cost, risk reduction, and the quoting/hedging interaction.
\end{enumerate}

\subsection{\texorpdfstring{Problem Set 9 --- Backtesting, Risk, and
Taker Flow \emph{(Python/PyTorch --- Milestone
3)}}{Problem Set 9 --- Backtesting, Risk, and Taker Flow (Python/PyTorch --- Milestone 3)}}\label{problem-set-9-backtesting-risk-and-taker-flow-pythonpytorch-milestone-3}

\begin{enumerate}
\def\labelenumi{\arabic{enumi}.}
\tightlist
\item
  \textbf{Event-driven backtest.} Implement the harness with
  deterministic replay; validate by running an identical config twice
  for bitwise results.
\item
  \textbf{Fill simulator.} Model fill probability as a function of quote
  distance and opposite-side depth; calibrate on a training month and
  validate on a disjoint holdout.
\item
  \textbf{Full backtest.} Run the multi-Greek engine on Coincall BTC
  options; report mean/std daily P\&L, Sharpe (with block-bootstrap CI),
  max drawdown, and aggregated-Greek time series.
\item
  \textbf{P\&L attribution.} Decompose daily P\&L into spread,
  inventory, gamma, and hedging; emit a CSV for a stacked-area plot;
  assert the streams sum to total P\&L.
\item
  \textbf{Drawdown forensic + taker analysis.} Identify the
  worst-drawdown day, report the environment, the positions that drove
  the loss, and a counterfactual. Using the guest lecture, characterize
  the taker flow on that day and how it affected your fills.
\end{enumerate}

\emph{Deliverable: Milestone 3 (see syllabus §2 for acceptance
criteria).}

\subsection{\texorpdfstring{Problem Set 10 --- Advanced Extension and
Final Report
\emph{(Python/PyTorch)}}{Problem Set 10 --- Advanced Extension and Final Report (Python/PyTorch)}}\label{problem-set-10-advanced-extension-and-final-report-pythonpytorch}

\begin{enumerate}
\def\labelenumi{\arabic{enumi}.}
\tightlist
\item
  \textbf{Proposal.} Submit a one-page proposal (problem, method,
  expected results, risks) for your chosen advanced Avellaneda--Stoikov
  extension by end of day one.
\item
  \textbf{Implementation.} Implement the extension in PyTorch on top of
  the milestone infrastructure (multi-asset, adverse-selection,
  signal-driven, discrete-inventory, or robust).
\item
  \textbf{Evaluation.} Quantify the improvement (or not) over the Week-9
  baseline using the same risk and P\&L metrics, over the same window.
\item
  \textbf{Final report.} ≤ 30 pages, following the lecture's structure,
  incorporating all milestone work.
\item
  \textbf{Final presentation.} Twenty minutes to the cohort and invited
  industry guests.
\end{enumerate}

\end{document}
